%%% Кандидатская диссертация — полный шаблон (все элементы)
%%% Тематика: Теория линейных кодов и криптография
%%% Компиляция: xelatex → biber → xelatex → xelatex

\documentclass[font=modern]{cmcmsuthesis-phd}

\author{Сидоров Алексей Михайлович}
\title{Произведение Шура--Адамара линейных кодов
       в анализе теоретико-кодовых криптосистем}
\thesissetup{department={Кафедра информационной безопасности}}
\supervisor{name={д-р~физ.-мат.~наук, проф. Чижов~И.В.}}
\specialty{code={1.2.3},name={теоретическая информатика, кибернетика}}

\usepackage{biblatex}
\addbibresource{refs.bib}

\begin{document}
\maketitle
\tableofcontents

% ================================================================
\chapter*{Введение}
\addcontentsline{toc}{chapter}{Введение}

Теория кодов, исправляющих ошибки, находит приложения
в криптографии~\cite{shannon1948,macwilliams1979}.

\textbf{Целью} работы является исследование свойств произведения
Шура--Адамара линейных кодов.

\textbf{Задачи}:
\begin{enumerate}[beginpenalty=10000]
    \item исследовать структуру произведения кодов;
    \item получить оценки кодового расстояния;
    \item разработать алгоритм вычисления спектра весов;
    \item проанализировать стойкость криптосистем.
\end{enumerate}

% ================================================================
\chapter{Сведения из теории линейных кодов}\label{ch:codes}

\section{Основные определения}

Пусть $\mathbb{F}_q$~--- конечное поле порядка $q$.
Элементы $\mathbb{F}^n_q$~--- векторы-строки
$\vect{v}=(v_1,\dotsc,v_n)$.
Матрица $\matr{M}\in\mathbb{F}^{k\times n}_q$.
Группа $\str{S}_n$, подстановка $\elt{\sigma}$, код $\code{C}$,
алгоритм $\algo{Decode}$, сложность $\compl{O}(n^2)$,
шифр $\crypto{E}$.

\begin{definition}\label{def:nkcode}
$[n]_q$-код $\code{C}$ называется $[n,k]_q$-кодом, если
$\dim\code{C}=k$.
\end{definition}

\begin{definition}\label{def:generator}
Матрица $\matr{G}\in\mathbb{F}^{k\times n}_q$, строки которой~---
базис кода $\code{C}$, называется \emph{порождающей}.
\end{definition}

\begin{theorem}\label{thm:dual-dim}
Если $\code{C}$~--- $[n,k]_q$-код, то $\code{C}^\perp$~---
$[n,n-k]_q$-код и $(\code{C}^\perp)^\perp=\code{C}$.
\end{theorem}

\begin{proof}
Следует из теории линейных пространств и определения
ортогонального дополнения.
\end{proof}

\begin{corollary}\label{cor:parity}
Если $\matr{G}=(\matr{I}_k\mid\matr{G}')$~--- систематическая,
то $\matr{H}=(-\matr{G}'^\top\mid\matr{I}_{n-k})$~--- проверочная.
\end{corollary}

\begin{lemma}\label{lem:weight}
$d(\code{C})=\min_{\vect{c}\neq\vect{0}}\operatorname{wt}(\vect{c})$.
\end{lemma}

\begin{proposition}\label{prop:gen-count}
Число порождающих матриц $[n,k]_q$-кода:
$(q^k-1)(q^k-q)\cdots(q^k-q^{k-1})$.
\end{proposition}

\begin{remark}\label{rem:systematic}
Не у каждого кода есть систематическая матрица в узком смысле.
\end{remark}

\begin{example}\label{ex:hamming}
Код Хемминга $[7,4,3]_2$ исправляет одну ошибку.
\end{example}

\begin{problem}\label{prob:selfdual}
Описать все самодуальные $[n,n/2]_2$-коды при $n\leq 24$.
\end{problem}

\section{Кодовое расстояние и спектр весов}

\begin{definition}\label{def:distance}
Минимальное расстояние: $d=\min_{\vect{c}\neq\vect{0}}\operatorname{wt}(\vect{c})$.
\end{definition}

\begin{definition}\label{def:spectrum}
Спектр весов: $\vect{A}(\code{C})=(A_0,A_1,\ldots,A_n)$,
$A_i=|\{\vect{c}\in\code{C}\mid\operatorname{wt}(\vect{c})=i\}|$.
\end{definition}

% ================================================================
\chapter{Формулы}\label{ch:formulas}

\section{Ненумерованные формулы}
Inline: $x\approx\sin x$ при $x\to 0$.
Display: \[(x_1+x_2)^2=x_1^2+2x_1x_2+x_2^2.\]
Непрерывная дробь:
\[\frac{1}{\sqrt{2}+\displaystyle\frac{1}{\sqrt{2}+\displaystyle\frac{1}{\sqrt{2}+\cdots}}}\]

\section{Нумерованные формулы}
\begin{equation}\label{eq:euler}e=\lim_{n\to\infty}\left(1+\frac{1}{n}\right)^{\!n}.\end{equation}
\begin{equation}\label{eq:basel}\lim_{n\to\infty}\sum_{k=1}^n\frac{1}{k^2}=\frac{\pi^2}{6}.\end{equation}
Ссылки: \eqref{eq:euler} и~\eqref{eq:basel}.
\begin{equation}\label{eq:long}\begin{multlined}1+2+3+\dots+50+\dots+\\+96+97+98+99+100=5050.\end{multlined}\end{equation}

\section{Многострочные формулы}
\[\begin{aligned}f_W&=\min\!\left(1,\max\!\left(0,\frac{W/W_{\max}}{W_{\text{crit}}}\right)\right),\\f_T&=\min\!\left(1,\max\!\left(0,\frac{T_s/T_{\text{melt}}}{T_{\text{crit}}}\right)\right).\end{aligned}\]

\[|x|=\left\{\begin{alignedat}{2}& &x,\quad&\text{если }x\geqslant 0;\\&-&x,\quad&\text{если }x<0.\end{alignedat}\right.\]

\[|x|=\begin{cases}\phantom{-}x,&\text{если }x\geqslant 0;\\-x,&\text{если }x<0.\end{cases}\]

\begin{equation}\label{eq:hadamard}\begin{alignedat}{2}
(\code{C}_1\star\code{C}_2)&=\langle\vect{c}_1\star\vect{c}_2\mid\vect{c}_1\in\code{C}_1,\;\vect{c}_2\in\code{C}_2\rangle,&\quad&\\
\dim(\code{C}_1\star\code{C}_2)&\leq\min(n,\;k_1 k_2).&&
\end{alignedat}\end{equation}

\section{Субформулы}
\begin{subequations}\label{eq:sub1}\begin{gather}y=x^2+1\label{eq:s1a}\\y=2x^2-x+1\label{eq:s1b}\\y=3x^2-4x+8\label{eq:s1c}\end{gather}\end{subequations}
\begin{subequations}\label{eq:sub2}\begin{align}y&=x^3+x^2+x+1\label{eq:s2a}\\y&=x^2\label{eq:s2b}\end{align}\end{subequations}
Ссылки: \eqref{eq:s1a}, \eqref{eq:s1b}, \eqref{eq:s2a}, \eqref{eq:s1c}.

Система (Лоренц): \[\left\{\begin{array}{rl}\dot x=&\sigma(y-x)\\\dot y=&x(r-z)-y\\\dot z=&xy-bz\end{array}\right..\]

Матрица: \[\left(\begin{array}{ccc}a_{11}&\ldots&a_{1n}\\\vdots&\ddots&\vdots\\a_{n1}&\ldots&a_{nn}\end{array}\right).\]

Длинная inline: $a_1+a_2+a_3+a_4+a_5+a_6+a_7+a_8+a_9+a_{10}+a_{11}+a_{12}+a_{13}+a_{14}+a_{15}+a_{16}+a_{17}+a_{18}+a_{19}+a_{20}=\sum_{i=1}^{20}a_i.$

\section{Греческие буквы и алфавиты}
Обычные: $\alpha\beta\gamma\delta\epsilon\theta\lambda\mu\pi\sigma\phi\omega\Gamma\Delta\Omega$.
Прямые: $\symup{\alpha}\symup{\beta}\symup{\gamma}\symup{\delta}\symup{\Gamma}\symup{\Omega}$.
Жирные: $\symbf{\alpha}\symbf{\beta}\symbf{\gamma}\symbf{\Gamma}\symbf{\Omega}$.
Жирные прямые: $\symbfup{\alpha}\symbfup{\beta}\symbfup{\gamma}\symbfup{\Gamma}\symbfup{\Omega}$.

\[\begin{aligned}\mathcal{ABCDEFGHIJKLMNOPQRSTUVWXYZ}\\\mathfrak{ABCDEFGHIJKLMNOPQRSTUVWXYZ}\\\mathbb{ABCDEFGHIJKLMNOPQRSTUVWXYZ}\end{aligned}\]

Русские операторы: $\tan\alpha$/$\cot\beta$/$\csc\gamma$; $\le$/$\leqslant$; $\phi$/$\varphi$; $\epsilon$/$\varepsilon$; $\kappa$/$\varkappa$; $\emptyset$/$\varnothing$.

Отступы:
\[\begin{aligned}\backslash!&\quad f(x)=x^2\!+3x\!+2\\\text{default}&\quad f(x)=x^2+3x+2\\\backslash,&\quad f(x)=x^2\,+3x\,+2\\\backslash;&\quad f(x)=x^2\;+3x\;+2\\\backslash\text{quad}&\quad f(x)=x^2\quad+3x\quad+2\end{aligned}\]

% ================================================================
\chapter{Форматирование, таблицы, изображения и списки}

\section{Форматирование текста}
\begin{description}
\item[Жирный:] \textbf{жирный}.\item[Курсив:] \emph{курсив}.
\item[Жирный курсив:] \textbf{\emph{жирный курсив}}.
\item[Подчёркивание:] \underline{подчёркнутый}.
\item[Зачёркивание:] \sout{зачёркнутый}.
\item[Sans:] \textsf{sans serif font}.
\end{description}

\section{Таблицы}
\begin{table}[htbp]\caption{Простая таблица}\label{tab:simple}\centering
\begin{tabular}{|l|c|c|c|}\hline Код&$n$&$k$&$d$\\\hline
Хемминга&7&4&3\\\hline Голея&23&12&7\\\hline РМ&32&16&8\\\hline
\end{tabular}\end{table}

\begin{table}[htbp]\caption{Booktabs}\label{tab:booktabs}\centering
\begin{tabular}{lrcc}\toprule
\textbf{Код}&\textbf{$|\code{C}|$}&\textbf{$d(\code{C}\star\code{C})$}&\textbf{Стойкость}\\\midrule
$[15,5,7]$&32&3&высокая\\
$[31,16,7]$&65536&3&средняя\\
$[63,32,11]$&$2^{32}$&5&высокая\\\bottomrule
\end{tabular}\end{table}

\begin{longtable}{@{}rr@{}}\caption{Длинная таблица}\label{tab:long}\\
\toprule $i$&$A_i$\\\midrule\endfirsthead
\toprule $i$&$A_i$\\\midrule\endhead\bottomrule\endfoot
0&1\\1&0\\2&0\\3&7\\4&7\\5&18\\6&18\\7&16\\\end{longtable}
Ссылка: таблица~\ref{tab:booktabs}.

\section{Изображения}
%% \begin{figure}[htbp]\centering\includegraphics[width=0.8\textwidth]{images/tanner.png}\caption{Граф Таннера кода}\label{fig:tanner}\end{figure}
%% \begin{figure}[htbp]\centering\includegraphics[scale=0.5]{images/spectrum.png}\caption{Спектр весов (масштаб 0.5)}\label{fig:spectrum}\end{figure}
%% Ссылка: рисунок~\ref{fig:tanner}.

\section{Списки}
\subsection{Нумерованный}
\begin{enumerate}\item Первый\begin{enumerate}\item Второй\begin{enumerate}\item Третий\begin{enumerate}\item Четвёртый\end{enumerate}\end{enumerate}\end{enumerate}\end{enumerate}
\subsection{Ненумерованный}
\begin{itemize}\item Первый\begin{itemize}\item Второй\begin{itemize}\item Третий\begin{itemize}\item Четвёртый\end{itemize}\end{itemize}\end{itemize}\end{itemize}

% ================================================================
\chapter{Алгоритмы и листинги}

\begin{algorithm}[H]\caption{Вычисление спектра весов}\label{alg:spectrum}
\KwInput{порождающая матрица $\matr{G}\in\mathbb{F}^{k\times n}_2$}
\KwOutput{спектр $\vect{A}=(A_0,\ldots,A_n)$}
$\vect{A}\gets(0,\ldots,0)$\;
\For{$\vect{a}\in\mathbb{F}^k_2$}{
    $\vect{c}\gets\vect{a}\matr{G}$\tcp*{кодовое слово}
    $w\gets\operatorname{wt}(\vect{c})$\;
    \Comment{обновляем счётчик}
    \uIf{$w=0$}{$A_0\gets A_0+1$\;}
    \uElseIf{$w=n$}{$A_n\gets A_n+1$\;}
    \Else{$A_w\gets A_w+1$\;}
}
\Return{$\vect{A}$}\;
\end{algorithm}
Ссылка: алгоритм~\ref{alg:spectrum}.

\begin{lstlisting}[language=Python,caption={Python},label={lst:py}]
import numpy as np
def weight_spectrum(G):
    """Спектр весов двоичного кода."""
    k, n = G.shape
    A = np.zeros(n + 1, dtype=int)
    for i in range(2**k):
        # Формируем информационный вектор
        a = np.array([(i >> j) & 1 for j in range(k)])
        c = a @ G % 2
        A[np.sum(c)] += 1
    return A
G = np.array([[1,0,1,1],[0,1,1,0]])
print("Спектр:", weight_spectrum(G))
\end{lstlisting}

\begin{lstlisting}[language=C,caption={C},label={lst:c}]
#include <stdio.h>
/* Вес Хемминга */
int hamming_weight(unsigned int x) {
    int w = 0;
    // Считаем единичные биты
    while (x) { w += x & 1; x >>= 1; }
    return w;
}
int main(void) {
    printf("wt(0b1011) = %d\n", hamming_weight(0xB));
    return 0;
}
\end{lstlisting}

\begin{lstlisting}[style=pseudocode,caption={Декодирование}]
(*@\lstInput@*): вектор $\vect{y}\in\mathbb{F}^n_2$, проверочная $\matr{H}$
(*@\lstOutput@*): ближайшее кодовое слово
procedure Decode($\vect{y}$, $\matr{H}$)
begin
    $\vect{s} \gets \matr{H}\vect{y}^\top$
    if $\vect{s} = \vect{0}$ then return $\vect{y}$
    for $j = 1$ to $n$ do
        if $\vect{s}$ совпадает со столбцом $j$ then
            return $\vect{y} \oplus \vect{e}_j$
    end
    return nil
end
\end{lstlisting}

\begin{lstpseudocode}[caption={Произведение Шура},label={alg:schur}]
(*@\lstInput@*): базисы $\code{C}_1$, $\code{C}_2$
(*@\lstOutput@*): базис $\code{C}_1 \star \code{C}_2$
function SchurProduct($\code{C}_1$, $\code{C}_2$)
begin
    $B \gets \emptyset$
    for each $\vect{c}_1 \in \text{basis}(\code{C}_1)$ do
        for each $\vect{c}_2 \in \text{basis}(\code{C}_2)$ do
            $B \gets B \cup \{\vect{c}_1 \star \vect{c}_2\}$
    end
    return RowReduce($B$)
end
\end{lstpseudocode}

% ================================================================
\chapter{Библиография и ссылки}
Одиночная~\cite{shannon1948}. Книга~\cite{knuth1984}.
Множественная~\cite{shannon1948,knuth1984,macwilliams1979}.
Ссылки: теорема~\ref{thm:dual-dim}, определение~\ref{def:nkcode},
формула~\eqref{eq:hadamard}, таблица~\ref{tab:booktabs},
алгоритм~\ref{alg:spectrum}.

% ================================================================
\chapter*{Заключение}
\addcontentsline{toc}{chapter}{Заключение}
\begin{enumerate}
\item Получены оценки расстояния произведения Шура--Адамара.
\item Разработан алгоритм вычисления спектра (алгоритм~\ref{alg:spectrum}).
\item Проведён анализ стойкости криптосистем.
\end{enumerate}

\clearpage\listoftables\addcontentsline{toc}{chapter}{Список таблиц}
\printbibliography[heading=bibintoc]

\appendix
\chapter{Доказательства вспомогательных лемм}
\section{Лемма о весе}
Доказательство леммы~\ref{lem:weight}.
\section{Оценка размерности}
Дополнительные выкладки.

\chapter{Результаты экспериментов}
Спектры весов для различных кодов.
\section{Коды БЧХ}
Таблицы.

\end{document}
